\NeedsTeXFormat{LaTeX2e}

% Rahmenumgebung f�r Studien Diplomarbeiten
% Erstellt von Oomke Weikert und Florian Keiler
% �nderungen im file _changelog.txt

% Titel und Autor der Arbeit unten im 
% \hypersetup Kommando �ndern
% ...und nat�rlich auf der Titelseite (title.tex bzw. title_en.tex)

% um pdf-file zu erzeugen:
% compilieren mit:
% pdflatex 
% bibtex
% pdflatex  
% pdflatex 
% (Hinweis: Das pdf-file darf bei Aufruf von pdflatex
% nicht im Adobe Reader ge�ffnet sein. Wenn man das
% pdf-file mit Ghostview �ffnet, muss es nicht geschlossen
% werden, und man kann dort die gerade bearbeitete Seite
% offen lassen)

% Im figures Ordner m�ssen die Bilder z.B. in pdf oder jpg Format liegen,
% mit pdflatex k�nnen *keine* eps Bilder benutzt werden.
% Das pdf file darf beim Kompilieren nicht 
% im Acrobat Reader ge�ffnet sein!!!

% um ps-file zu erzeugen:
% compilieren mit:
% latex 
% bibtex
% latex 
% latex 
% dvips

% dvips Aufruf f�r Type-1 Schriften: 
% dvips -t a4 -Ppdf 
% (alte GhostScript-Version: dvips -t a4 -Ppdf -G0) 
% Umwandeln in pdf mit Acrobat Distiller 
% oder mit ps2pdf (in GhostScript enthalten)


\documentclass[a4paper, twoside, 11pt]{report}

\newif\ifmakeindex
\makeindextrue % generate index
%\makeindexfalse % don't generate index

\ifmakeindex
% fuer Stichwortverzeichnis
\usepackage{makeidx}

% Stichwortverzeichnis erstellen
\makeindex
\fi

\newif\ifenglish
\englishfalse  % german document

\usepackage[margin=1cm,format=hang,font=small,labelfont=bf,textfont=sl]{caption}
% angepasste Bildunterschriften
% Doku siehe beigef�gtes pdf-file

\usepackage{subfigure}
% erm�glicht subfigures aufzubauen
% Doku ftp://tug.ctan.org/pub/tex-archive/obsolete/macros/latex/contrib/subfigure/subfigure.pdf

%	\usepackage{ngerman}
% Trennmuster f�r neue deutsche Rechtschreibung
\usepackage[ngerman]{babel}


\usepackage[latin1]{inputenc}
% Unterstuetzen von deutschen Umlauten

\usepackage{t1enc}
% Verwenden von DC-Fonts (erlaubt richtiges Trennen
% auch in Worten mit Umlauten)

%\usepackage{times} 
  % Saves a lot of ouptut space in PDF...
  % ..after conversion with ps2pdf or Adobe Distiller	 
	% �ndert aber die Schriftart! (Times statt computer modern roman)
\usepackage{courier} % fett *und* typewriter nur damit m�glich?
% �ndert normale Schrift nicht!

\usepackage{bibgerm}
% Deutsche Bezeichnungen


\usepackage{cite}
% Automatisches Zusammenfassen von Literaturstellen

\usepackage{fancyheadings}	
%STY-file fancyheadings.sty nicht standardm��ig in miktex enthalten
%evtl. durch fancyhdr ersetzen

\usepackage{float}
%This package improves the interface for defining floating objects such
%as figures and tables in LaTeX.  It adds the notion of a `float style'
%that governs appearance of floats.  New kinds of floats may be defined
%using a \newfloat command analogous to \newtheorem.  This style option
%also incorporates the functionality of David Carlisle's style option
%`here', giving floating environments a [H] option which means `PUT IT
%HERE' (as opposed to the standard [h] option which means `You may put
%it here if you like').

\usepackage{verbatim}		
%This package reimplements the L A T E X verbatim and verbatim* envi- 
%ronments. In addition it provides a comment environment that skips any 
%commands or text between \begin{comment} and the next \end{comment}. 
%It also defines the command verbatiminput to input a whole file verbatim. 

\usepackage{amsmath}
\usepackage{amsfonts}
% erm�glicht Mengennotationen, wie z.B. f�r Nat�rliche Zahlen

%\usepackage{rotating}

\usepackage{array}
% wird benutzt von macro.tex f�r \Case

\usepackage{ifthen}
% wird benutzt von macro.tex \ifthenelse

\usepackage{longtable}
% f�r lange Tabellen gr��er als eine Seite mit \begin{longtable} .. \end{longtable}

\usepackage{varioref}
% Kommando \vref verweist mit Seitenzahl

\usepackage{readfile}
% ASCII-File einbinden mit
% \readit{fftdb.m}{\tt} 
% 2. Argument (\tt) gibt zu benutzende Schriftart an 
\usepackage{shapepar}
\usepackage{setspace}
\usepackage{colortbl}
\usepackage{color}
\usepackage{xcolor} 

\usepackage{listings}

\lstloadlanguages{}

\lstloadlanguages{C,C++,Java,Matlab,HTML,TeX,XML}


%\lstset{}% restore default
%# title={Titel}:  gibt einen Titel zur Umgebung an (erscheint �ber dem Code zentriert)
%# caption={Titel}:  
%# label=Name:

\lstset{
				frame=lines, %Umrandung (single|none|shadowbox|lines|bottomline|topline|leftline)
				framerule=1pt, %Rahmenbreite
				tabsize=4, %Anzahl der Zeichen f�r ein TAB
				backgroundcolor=\color{lightgray}, %Hintergrundfarbe
				emph={}, %hebt die angegebenen W�rter hervor
				emphstyle=\underbar, %unterstreicht hervorgehobene W�rter
				columns=fixed, %Zeichenabstand (fixed | flexible | fullflexible)
				lineskip=0pt, %Zeilenabstand
				basicstyle=\ttfamily,%\ttfamily,
				identifierstyle=\color{black},
        commentstyle=\color{darkgreen},
        stringstyle=\color{viola},
        keywordstyle=\color{darkblue},
				ndkeywordstyle=\color{black},
				showspaces=false,
				showtabs=false,
				numbers=none, %Zeilennummern (none|left|right)
				%numbertype=\ttfamily,
				breaklines=true,
        captionpos=b,
        extendedchars=false
}

%\usepackage{fancybox}
%\usepackage{theorem}
%\usepackage{amsbsy}
%\usepackage{nomencl}        % unterst�tzt Symbolverzeichnisse
%\usepackage{makeidx}
%\usepackage{multind}

\def\boxes{yes}
% \boxes == yes makes boxes around desired formulas by \Mbox
% wird benutzt von macros.tex
 

% pdf-tex settings:
% ------------------------
% detect automatically if run by latex or pdflatex
\newif\ifpdfout
\ifx\pdfoutput\undefined
   \pdfoutfalse
\else
   \pdfoutput=1
   \pdfouttrue
\fi

\ifpdfout % compiling with pdflatex
   \usepackage[pdftex]{graphicx}
   \DeclareGraphicsExtensions{.pdf, .png, .jpg}
   \usepackage[pdftex,
	bookmarks,
	%colorlinks=false, % instead of colors, now boxes are used for links
	colorlinks=true,
	urlcolor=black, %blue,
	linkcolor=black, %red, %normal internal links
	citecolor=black, %green, %citation links
	%pagebackref, %link from references back to page of citation
	linktocpage, 
	% im Inhaltsverzeichnis Link auf Seitenzahl (sonst Probleme bei langen Zeilen)
	%breaklinks = true, % f�r Links l�nger als 1 Zeile
	%hypertexnames = false,  % f�r Links zu Figures?
	bookmarksopen, %open all bookmark folders
	bookmarksnumbered, %use section numbers with bookmarks
    pdfpagemode=UseOutlines, %show bookmarks
	% http://www.tex.ac.uk/cgi-bin/texfaq2html?label=pdfpagelabels
	plainpages=false, % eigene Seitenanker f�r r�mische/arabische Seitenzahlen
	pdfpagelabels, % im Abode Reader Seitenzahl als z.B. "iii (3 von 20)" anzeigen
%	pdfstartview=FitH
	pdfstartview=FitV
	]{hyperref}
    \pdfadjustspacing=1                %%% force LaTeX-like character spacing
    \pdfcompresslevel=9
	\pdfcatalog{                 
	% Catalog dictionary of PDF output. 
    % /PageMode /UseNone          
    % /URI (http://www.fi.muni.cz/)
	%       
	% pdfscreen-like setting might look like:
	%     /PageMode /none 
	%     /ViewerPreferences << 
	%         /HideToolbar true            
	%         /HideMenubar true 
	%         /HideWindowUI true 
	%         /FitWindow true 
	%         /CenterWindow true 
	%	
	% /PageMode determines how Acrobat displays the document on startup. 
	% The possibilities for the latter are explained below:
	% Supported /PageMode values.
	% /UseNone 		neither outline nor thumbnails visible
	% /UseOutlines 	outline visible
	% /UseThumbs 	thumbnails visible
	% /FullScreen 	full--screen mode
	% In full--screen mode, there is no menu bar, window controls, 
	% nor any other window present. The default setting is /UseNone.
	} %end of \pdfcatalog
\else % compiling with latex
  \usepackage[dvips]{graphicx}
%  \usepackage{color}
  \DeclareGraphicsExtensions{.eps}
  \usepackage[
    dvips,
%	ps2pdf,  % statt dvips, Unterschied???
	bookmarks,
  	colorlinks=false,  % for final paper without colors
%  	colorlinks=true,
	linktocpage, 	
	% im Inhaltsverzeichnis Link auf Seitenzahl (sonst Probleme bei langen Zeilen)
	%breaklinks = true, % f�r Links l�nger als 1 Zeile
	% funktioniert *nicht* korrekt f�r lange Zeilen mit gs < 7.05.3 ?
%	pdfstartview=FitH, % funktioniert NICHT mit Adobe Distiller
	bookmarksopen, %open all bookmark folders
	bookmarksnumbered, %use section numbers with bookmarks
    pdfpagemode=UseOutlines, %show bookmarks
	pdfstartview=FitV
   ]{hyperref}
  % hyperrefs are active is the pdf file after conversion
  \hypersetup{
	pdfcreator  = {LaTeX with hyperref package},
	pdfproducer = {dvips + ps2pdf}
  }
\fi

% HIER TITEL UND AUTOR ANPASSEN!!!
\hypersetup{
  pdftitle={Studien/Diplomarbeit},
	pdfauthor={Autor},
	pdfsubject  = {Studien/Diplomarbeit, UniBwH, Professur ANT, \today},
	pdfkeywords = {}
}

% ------------------------


%Breite der Bild/Tabellenbeschriftungen:
%\setlength{\LTcapwidth}{\captionwidth}

%Schriftgr��e/Zeilenvorschub in Tabellen:
\def\tablefontsize{\normalsize}
\def\tablespacing{1.0} % linespacing in tables
  
\newcommand{\bc}{\begin{center}}
\newcommand{\ec}{\end{center}} 

\newcommand{\be}{\begin{equation}}
\newcommand{\ee}{\end{equation}} 
\newcommand{\bea}{\begin{eqnarray}}
\newcommand{\eea}{\end{eqnarray}}

\newcommand{\bi}{\begin{itemize}}
\newcommand{\ei}{\end{itemize}}

\newcommand{\bmp}{\begin{minipage}[c]{\columnwidth} \bc}
% argument: vertical position: c, t, b
\newcommand{\emp}{\ec \end{minipage} }

\newcommand{\beg}{\hspace{0.5cm}\begin{minipage}{14cm}}
\newcommand{\eeg}{\end{minipage}\vspace{1cm}}


%%%%%%%%%%%%%% tables %%%%%%%%%%%%%%%%%%%%%%%
\newcommand{\bt}[1]
% #1: table placing: htbp
    {    \begin{table}[#1] %htbp 
         \begin{center}
         \renewcommand{\baselinestretch}{\tablespacing}
         \tablefontsize
    }

%Tabelle mit gr��erem Zeilenvorschub:
\newcommand{\btline}[2]
% #1: table placing: htbp
% #2: linespacing in table, e.g. 1.4
{   \def\tablespacing{#2}
    \bt{#1}
} 

\newcommand {\et}[2] 
% #1: caption
% #2: label without 'tab:'
    {
     \caption{#1.} 
     \label{tab:#2}
     \end{center} 
     \end{table}
     \def\tablespacing{1.0} 
    % reset linespacing for next table
}

\newcommand {\btab}{\begin{tabular}}
\newcommand {\etab} {\end{tabular}}

\newcommand{\mc}{\multicolumn}
\newcommand{\mr}{\multirow}

% columns in math mode:
\newcolumntype{C}{>{$}c<{$}}
\newcolumntype{L}{>{$}l<{$}}
\newcolumntype{R}{>{$}r<{$}}

\newcommand{\pbs}[1]{\let\temp=\\#1\let\\=\temp}
%preserve backslash, s. companion, p108



%%%%%%%%%%%%%% figures %%%%%%%%%%%%%%%%%%%%%%%
\newcommand{\fig}[4]{
% Bild mit 80% der Seitenbreite
% arguments:	
% #1: file without extension eps
% #2: caption
% #3: label
% #4: placing of the figure: e.g. htbp
\begin{figure}[#4]			      
\begin{center}
   \includegraphics[width=.8\columnwidth]{figures/#1}
   \caption{#2.}
   \label{fig:#3}
\end{center}
\end{figure}
}

\newcommand{\figscale}[5]{
% arguments:	
% #1: file without extension eps
% #2: caption
% #3: label
% #4: placing of the figure: e.g. htbp
% #5: width of figure
\begin{figure}[#4]			      
\begin{center}
   \includegraphics[width=#5\columnwidth]{figures/#1}
   \caption{#2.}
   \label{fig:#3}
\end{center}
\end{figure}
}

\newcommand{\figscaletwo}[6]{
% arguments:	
% #1: file1 without extension eps
% #2: file2 without extension eps
% #3: caption
% #4: label
% #5: placing of the figure: e.g. htbp
% #6: width of figure
\begin{figure}[#5]			      
\begin{center}
   \includegraphics[width=#6\columnwidth]{figures/#1}~\\
   \includegraphics[width=#6\columnwidth]{figures/#2}
   \caption{#3.}
   \label{fig:#4}
\end{center}
\end{figure}
}

\newcommand{\figtwo}[5]{
% two figures labeled (a) and (b) side by side
% arguments:	
% #1: file1 without extension eps
% #2: file2 without extension eps
% #3: caption without fullstop
% #4: label
% #5: placing of the figure: e.g. htbp
\begin{figure}[#5]			      
\begin{center}
    \btab{cc} 
   \includegraphics[width=0.45\columnwidth]{figures/#1}
    &	 
   \includegraphics[width=0.45\columnwidth]{figures/#2}
    (a)&(b)
    \etab 
	\caption{#3.}
	\label{fig:#4}
\end{center}
\end{figure}
}

\newcommand{\figTwo}[6]{
% two figures labeled (a) and (b) on top/bottom
% arguments:	
% #1: file1 without extension eps
% #2: file2 without extension eps
% #3: caption without fullstop
% #4: label
% #5: placing of the figure: e.g. htbp
% #6: width of figures as ratio, e.g. 0.7
\begin{figure}[#5]			      
\begin{center}
   \includegraphics[width=#6\columnwidth]{figures/#1}
(a)\\
\vspace{1 em}
   \includegraphics[width=#6\columnwidth]{figures/#2}
    (b)\\
    \caption{#3.}
    \label{fig:#4}
\end{center}
\end{figure}
}

%%%%%%%%%%%%%%%%%%%%%%%%%%%%%%%%%%%%%%%%%%%%%%

\newcommand{\C}[1]{\texttt{#1}} %Typewriter-Font
\newcommand{\tild}{\~~\hspace{-1 ex}} %Tilde ~

\newcommand{\equi}{\Leftrightarrow} %�quivalenz <=>
\newcommand{\concl}{\Rightarrow}    %daraus folgt =>
%Matrix:
\newcommand{\matr}[2]{\left ( \begin{array}{#1} #2 \end{array} \right )}
% #1 alignment of colums, e.g. ccc
% #2 contents of matrix
%Vektor:
\newcommand{\vect}[1]{\left ( \begin{array}{c} #1 \end{array} \right )}
%Fettschrift im Mathemodus:
\newcommand{\mbf}[1]{\mathbf{#1}} %boldface
\newcommand{\gbf}[1]{\boldsymbol{#1}} %boldface greek letters
%Verweis auf Gleichungen mit Klammern um Gl.-Nr.
%\newcommand{\eqref}[1]{(\ref{#1})}
% \eqref defined in package amsmath


\newcommand{\freqorig}{\ensuremath{\circ\hspace{-.1em}\mbox{---}\!\bullet}}
\newcommand{\freq}{\mbox{~}{\circ\hspace{-.1em}\mbox{---}\!\bullet}\mbox{~}}
	% correspondence symbol time<->freq 
\newcommand{\freqswap}{\ensuremath{\bullet\!\mbox{---}\hspace{-.1em}\circ}}
\newcommand{\freqv}{\begin{turn}{90} \freqswap \end{turn}}
	% vertical correspondence symbol
	% time 
	% freq
\newcommand{\timeh}{\mbox{~}{\bullet\!\mbox{---}\hspace{-.1em}\circ}\mbox{~}}
	% correspondence symbol freq<->time
\newcommand{\timev}{\begin{turn}{90} \freqorig \end{turn}}
	% vertical correspondence symbol
	% freq
	% time

\newcommand{\Case}[1]{\left\{ \begin{array}{ll} #1
    \end{array}\right. } 


\newcommand{\Mbox}[1]{
\ifthenelse{\equal{\boxes}{yes}}
{%\ensuremath
\fbox{$\displaystyle #1$}}
{#1}}

%box around a equationarray
%argument is the contents of the eqnarray environment
\newcommand{\Mboxarray}[1]{
\ifthenelse{\equal{\boxes}{yes}}
{\bc\fbox{\begin{Beqnarray} #1 \end{Beqnarray}}\ec}
{\bea #1 \eea}
}


\newcommand{\suml}{\sum\limits}
% limits of sum UNDER the sign instead of beside it, e.g. in fractions

\newcommand{\intl}{\int\limits}
% limits of int UNDER the sign instead of beside it


\newcommand{\logd}{\mbox{ld\,}} %base 2 logarithm

\newcommand{\erw}[1]{E\left\{#1\right\}} %Erwartungswert


% make ',' an ordinary Symbol in decimal numbers
% Unterdr�ckung des Zwischenraums hinter einem Komma (z.B. in Dezimalzahlen)
% Steht im Mathemodus hinter dem Komma ein Leerzeichen, wird es als Trennzeichen
% (mit Zwischenraum) benutzt, sonst als Dezimalkomma.
 \mathchardef\CommaOrdinary="013B
 \mathchardef\CommaPunct   ="613B
 \mathcode`,="8000   % , im Math-Mode aktiv ("8000) machen
 {\catcode`\,=\active
  \gdef ,{\obeyspaces\futurelet\next\CommaCheck}}
 \def\CommaCheck{\if\space\next\CommaPunct\else\CommaOrdinary\fi}

% define the size of the region where text can appear on the page
\setlength{\hoffset}{-1in}
\textwidth 15.7cm
\textheight 24cm %24.4cm
\topmargin -1cm %-1.6cm
\oddsidemargin 2.7cm
\evensidemargin 2.6cm
%\oddsidemargin -3cm
%\evensidemargin -3cm
\headheight 15pt%[12]
\headsep 1.2cm%[25pt]
%\setlength{\footheight}{15pt} %[12]
%\columnsep 0.8cm

% setzt den Einzug bei neuem Absatz zu null
\setlength{\parindent}{0pt}
% oder \parindent0cm

%Abstand vor neuem Absatz
%wird in abstract.tex benutzt
\def\parspacing{2.5ex}
\setlength{\parskip}{\parspacing} 

% enable following command if you don't want extra space after points
\frenchspacing

\renewcommand{\textfraction}{0.1}
% Mindestanteil an Text auf einer Seite, 
% in Prozent, default: 0.2
 
\renewcommand{\topfraction}{1}
% Maximaler Anteil der von Gleitobjekten am Kopf der Seite 
% eingenommen werden kann, in Prozent, default: 0.7

\setcounter{topnumber}{10} 
% number of float objects at top of page, default: 2

\setcounter{totalnumber}{10} 
% number of float objects on one page, default: 3


\setcounter{secnumdepth}{3}
% depth of numbering, see companion p.20-21
% for report:   chapter: 0, section: 1, subsection: 2
%       subsubsection:3,paragraph: 4,   subparagraph: 5

\setcounter{tocdepth}{2}

\setlongtables
% for equal column width on different pages of a longtable, see companion, p 122
% has to be disabled to reset, only widest width is saved, 
%so there's an error, if this entry is shortened!


% Umdefinition der Bezeichner
\addto\captionsngerman{% Eigene Bezeichner: Abb. -> Bild, Tab.-> Tabelle
	\renewcommand{\figurename}{Bild}%
	\renewcommand{\tablename}{Tabelle}%
	\renewcommand{\lstlistlistingname}{Quellcodeverzeichnis} % default: "Listings"
	\renewcommand{\lstlistingname}{Quellcode} % default: "Listing"
}
\addto\captionsenglish{% 
	\renewcommand{\lstlistlistingname}{List of Sourcecodes} % default: "Listings"
	\renewcommand{\lstlistingname}{Sourcecode} % default: "Listing"
}


%\makeglossary

%\includeonly{title,abstract,erklaerung,vorwort,fanc_h_2,chap3,bib}

\begin{document}
\definecolor{gold}{rgb}{0.85,.66,0}
\definecolor{viola}{rgb}{0.5,0,0.5}
\definecolor{darkblue}{rgb}{0,0,.6}
\definecolor{darkred}{rgb}{.6,0,0}
\definecolor{darkgreen}{rgb}{0,.6,0}
\definecolor{red}{rgb}{.98,0,0}
\definecolor{lightgray}{rgb}{.99,.99,.97}


%HSU Corporate Design
\definecolor{hsurot}{cmyk}{0,1.00,0.43,0.185}
\definecolor{hsugrau}{cmyk}{0,0.085,0.15,0.43}
\definecolor{hsugelb}{cmyk}{0,0.15,0.60,0}
\definecolor{hsublau}{cmyk}{1.00,0.38,0,0.69}
\definecolor{hsugruen}{cmyk}{1.00,0,0.47,0.47}
\definecolor{hsuoliv}{cmyk}{0,0,1.00,0.47}
\definecolor{hsuorange}{cmyk}{0,0.43,1.00,0.18}
\definecolor{hsuhellrot}{cmyk}{0,0.76,0.83,0.11}

 

% 60%
\colorlet{hsurot60}{hsurot!60}
\colorlet{hsugrau60}{hsugrau!60}
\colorlet{hsugelb60}{hsugelb!60}
\colorlet{hsublau60}{hsublau!60}
\colorlet{hsugruen60}{hsugruen!60}
\colorlet{hsuoliv60}{hsuoliv!60}
\colorlet{hsuorange60}{hsuorange!60}
\colorlet{hsuhellrot60}{hsuhellrot!60}

% 30%
\colorlet{hsurot30}{hsurot!30}
\colorlet{hsugrau30}{hsugrau!30}
\colorlet{hsugelb30}{hsugelb!30}
\colorlet{hsublau30}{hsublau!30}
\colorlet{hsugruen30}{hsugruen!30}
\colorlet{hsuoliv30}{hsuoliv!30}
\colorlet{hsuorange30}{hsuorange!30}
\colorlet{hsuhellrot30}{hsuhellrot!30} % several predefined colors - including the HSU corporate Design.
\pagenumbering{roman} % sonst gibt es 2 mal Seiten 1 und 2 (arabisch)
% so werden Probleme bei hyperlinks vermieden

%\def\voff{1cm}
%\addtolength{\voffset}{-\voff}

%\vspace{-2cm}
\pdfbookmark[1]{Titelseite}{sec:title}  % Bookmark im pdf file
\begin{titlepage}
\label{sec:title}

%\enlargethispage{\voff}
\large 
\begin{center}

% Unilogo mit Link auf Webseite
\href{http://www.unibw-hamburg.de} 
{%\includegraphics[width = 9cm]{figures/Uni_farbe}
\includegraphics[width = 9cm]{figures/HSU-logo_farbe} 
}

\vspace{2cm}

\begin{minipage}{.9\linewidth}
{\centering\Huge\bf Titel der Studien-/Diplomarbeit mit einer l\"angeren Beschreibung des Themas\par} 
% \par ist n�tig f�r korrekten Zeilenabstand im Titel
\end{minipage}

\large 

%\vspace{1.5cm}
%{von\\[.3cm] {\bf \LARGE Autor}\\[1cm] 
%\huge	{ \bf \it Diplomarbeit} \\[1cm]}

\vspace{1.5cm}
{\huge{\bf \it Diplomarbeit} %\\[1cm]
}

\vspace{0.6cm}
{
\Large von \\[.5cm] 
{\bf \LARGE \huge Autor}\\[1.5cm]
}

\vfill
\begin{tabular}{ll}
Ausgabetermin:& xx. XXXXXX 200x\\
Abgabetermin:& xx. XXXXXX 200x\\
Betreuer:&  Dipl.-Ing. zzz\\
Hochschullehrer:&  Prof.~Dr.-Ing.~habil.~Udo Z\"olzer\\
\end{tabular}
\end{center}
\vspace{2cm}

\end{titlepage}

%\addtolength{\voffset}{\voff}


\thispagestyle{empty} % Titelr�ckseite (S. 2) ohne Seitennummer

\cleardoublepage

\newcounter{pageno}
\setcounter{pageno}{1} %titlepage = page 1, but pagenumber not printed
\addtocounter{pageno}{1}

\cleardoublepage
\pdfbookmark[1]{\abstractname}{sec:abstract}  % Bookmark im pdf file
\begin{abstract}
\label{sec:abstract}
\addtocounter{pageno}{1}
\setcounter{page}{\arabic{pageno}}
\thispagestyle{plain}
\begin{center}
\begin{minipage}[t]{0.8\linewidth} 
\setlength{\parskip}{\parspacing}

Dies ist ein Text, der nur ein Beispiel darstellt. Dies ist ein Text, der nur ein Beispiel darstellt. 

Dies ist ein Text, der nur ein Beispiel darstellt.

\end{minipage}
\end{center}

\end{abstract}



\cleardoublepage
\begin{center}
\pdfbookmark[1]{Erkl�rung}{sec:erklaer}  % Bookmark im pdf file
{\Huge \bf Erkl�rung}
\label{sec:erklaer}
\end{center}

Hiermit erkl\"are ich, dass die vorliegende Arbeit von mir selbst\"andig und nur unter Verwendung der angegebenen Quellen und 
Hilfsmittel erstellt wurde.

\vspace{2cm}

Hamburg, den xxx



\cleardoublepage
\pdfbookmark[1]{Vorwort}{sec:vorwort}  % Bookmark im pdf file
\chapter*{Vorwort}
\label{sec:vorwort}
bla bla bla

bla bla bla

\vspace{2cm}
Hamburg, xxx



\cleardoublepage
\pagestyle{fancyplain}

% Font of the header
\def\headfont{\bfseries}%\sffamily}

% Gew�nschte Formatierung der Kopfzeile, default: alles in Gro�buchstaben
\renewcommand{\chaptermark}[1]
{\markboth{\chaptername { }\thechapter. #1} 
{\chaptername { }\thechapter. #1}} 
% p 99
% without name ``Chapter'', ``Appendix'' etc:
%\renewcommand{\chaptermark}[1]{\markboth{\thechapter. #1}{\thechapter. #1}} % p 99


% for twoside option:
\renewcommand{\chaptermark}[1]{\markboth{\chaptername { }\thechapter. #1}{}}
%\chaptername { }\thechapter. #1}}
\lhead[\fancyplain{\headfont\thepage}{\headfont\thepage}]
{\fancyplain{\headfont\rightmark}{\headfont\rightmark}}

\renewcommand{\sectionmark}[1]{\markright{\thesection\ #1}}

%\lhead{\fancyplain{\headfont \rightmark}{\headfont\rightmark}}
\rhead[\fancyplain{\headfont\leftmark}{\headfont\leftmark}]
{\fancyplain{\headfont\thepage}{\headfont\thepage}}
%\lfoot{\tt Arbeit \emph{Titel} -- Autor}
%\rfoot{\tt Stand \today} %aktuelles Datum
\cfoot{}
%\setlength{\plainfootrulewidth}{0pt}
%\renewcommand{\footrulewidth}{0pt}
\setlength{\plainheadrulewidth}{0.4pt}
 % set fancy headings
%$$$$$$$$$$$$$$$ tableofcontents $$$$$$$$$$$$$$$$$
\ifenglish
	\newcommand{\chapname}{Chapter}
\else
	\newcommand{\chapname}{Kapitel}
\fi % end of \ifenglish\renewcommand{\chaptername}{} %if name appears in header

\setlength{\parskip}{0ex}
\renewcommand{\baselinestretch}{1}
\normalsize

\cleardoublepage % to get correct page no. in TOC
\pdfbookmark[1]{\contentsname}{toc}  % Bookmark auf Inhaltsverzeichnis im pdf file
\tableofcontents


\setlength{\parskip}{1ex}
\cleardoublepage % to get correct page no. in TOC
\phantomsection
\setcounter{lofdepth}{2}
% subfigures werden mit aufgelistet in der LOF
\addcontentsline{toc}{chapter}{\protect\numberline{\listfigurename}}
\listoffigures

\cleardoublepage % to get correct page no. in TOC
\phantomsection
\addcontentsline{toc}{chapter}{\protect\numberline{\listtablename}}
\listoftables

\cleardoublepage % to get correct page no. in TOC
\phantomsection
\addcontentsline{toc}{chapter}{\protect\numberline{\lstlistlistingname}}
\lstlistoflistings


\renewcommand{\baselinestretch}{1.0}
\normalsize

\cleardoublepage % to get correct page no. in TOC

\setlength{\parskip}{\parspacing}

\clearpage
\renewcommand{\chaptername}{\chapname} %if name appears in header

\cleardoublepage

% Hier wird das Symbolzeichnis eingef�gt
% Falls es nicht gew�nscht ist, sind die folgenden sechs Zeilen auszukommentieren
\cleardoublepage % to get correct page no. in TOC
\phantomsection
\markboth{\bfseries SYMBOLVERZEICHNIS}{\bfseries SYMBOLVERZEICHNIS}
\addcontentsline{toc}{chapter}{\protect\numberline{Symbolverzeichnis}}
\chapter*{Symbolverzeichnis}
\label{sec:symbolverzeichnis}


$x(n)=x'(n)+ j\cdot x''(n)$\\
$x'(n)$				Realteil\\
$x''(n)$			Imagin�rteil\\

\begin{table}[htbp]
\begin{tabular}[t]{ll}
HNV					&	Haupt- zu Nebenmaximumverh�ltnis.\\
MF					&	Merit-Faktor.\\
$\gamma$		&	Phasenoffsets des Kanals.\\
$\lambda$		&	Anzahl der Symbole pro Rahmen.\\
$\xi$				&	Anzahl der Pr�ambelsymbole.\\
$\psi$			&	Pilotsymbolanzahl.\\
$\alpha$		&	Zeitkompressionsfaktor.\\
cp					&	Clock-Precision in Simulink.\\
$r_{O}$			&	Output-Rate.\\
$r_{I}$			&	Input-Rate.\\
$\tau$			&	Anzahl an Pilotsymbolen pro Sequenz.\\
$S_{i}$			&	i-te - Datensequenz.\\
$\sigma$		&	Anzahl an Datensequenzen pro Rahmen.\\
$\beta_{R}$	&	Anstiegsfaktor f�r die Einh�llendenberechnung.\\
X						& Einschaltgrenze f�r die Rahmenerkennung.\\
$\beta_{F}$	&	Abfallfaktor f�r die Einh�llendenberechnung..\\
$t_{F}$			& Abfallzeit in Samples f�r die Einh�llendenberechnung.\\
$E_{S}$			& mittlere Symbolenergie\\		
K						&	Schwellwertfaktor f�r die Rahmenerkennung.\\
$\overline{Corr}$ &	Mittelwert des Korrelationsergebnisses.\\
$\sigma_{Corr}$	& Standardabweichung des Korrelationsergebnisses.\\
$\overline{Max}$ &	Mittelwert der Maxima des Korrelationssignals.\\
$\sigma_{Max}$ &	Standardabweichung der Maxima des Korrelationssignals..\\
S						&	dynamische Schwelle f�r die Rahmenerkennung.\\
$P_{F}$			&	Falschdetektionswahrscheinlichkeit bei der Rahmenerkennung.\\
$P_{N,S}$		&	Nichterkennungswahrscheinlichkeit bei der Rahmenerkennung.\\
$P_{E}$			&	Falscheinschaltwahrscheinlichkeit f�r die Rahmenerkennung.\\
$P_{Ab}$		&	Falschabschaltwahrscheinlichkeit f�r die Rahmenerkennung.\\
$E_{R}$			& Erwartungswert der Rayleighverteilung
\end{tabular}  
%\label{}  
\end{table}


\cleardoublepage

\setcounter{page}{1}
\renewcommand{\thepage}{\arabic{page}}

%============================================================================
%============================================================================

\cleardoublepage
\chapter{Erstes Kapitel}
\ifenglish
bla bla bla bla bla bla bla bla bla bla bla bla bla bla bla
``bla bla bla'' bla bla bla bla bla bla bla bla bla bla bla bla
bla bla bla bla bla bla bla bla bla bla bla bla bla bla bla
\else
bla bla bla bla bla bla bla bla bla bla bla bla bla bla bla
"`bla bla bla"' bla bla bla bla bla bla bla bla bla bla bla bla
\glqq bla bla bla\grqq~bla bla bla bla bla bla bla bla bla bla bla bla
\fi

bla bla bla bla bla bla bla bla bla bla bla bla bla bla bla
bla bla bla bla bla bla bla bla bla bla bla bla bla bla bla
bla bla bla bla bla bla bla bla bla bla bla bla bla bla bla
bla bla bla bla bla bla bla bla bla bla bla bla bla bla bla

\section{Abschnitt eins\index{Abschnitt eins}}
bla bla bla bla bla bla bla bla bla bla bla bla bla bla bla
bla bla bla bla bla bla bla bla bla bla bla bla bla bla bla
bla bla bla bla bla bla bla bla bla bla bla bla bla bla bla
bla bla bla bla bla bla bla bla bla bla bla bla bla bla bla

bla bla bla.
Ein Zitat: \cite{fli:mult} und \cite{Bauch_MAS}

Ein Equationarray:
\bea
\hat X(z)&=&P(z)[E(z)+\hat X(z)]\\
\equi\quad\hat X(z)&=&\frac{P(z)}{1-P(z)}E(z)\\
\Omega &=& 0,25 \pi\\
\mathbf{Ax} &=&\mathbf{a}
\eea

\index{�bertragungstest(1)} 
\index{Ubertragungstest@�bertragungstest(2)} % => Sortierung unter U
\index{"Ubertragungstest(3)} 
	% => Sortierung unter Ue bei Option -g -s index_umlauts.ist 
\index{\"Ubertragungstest(4)} 
\index{Ubers}
\index{Uebers}

Eine Gleichung
\be
\sum_{k=1}^{m}a^{(m)}_k r_{xx}(i-k)= r_{xx}(i)\quad,\:i=1,\ldots,m
\label{eq:normalgl}
\ee
und mit Box:
\be
\Mbox{
\sum_{k=1}^{m}a^{(m)}_k r_{xx}(i-k)= r_{xx}(i)\quad,\:i=1,\ldots,m
}.
\label{eq:normalglbox}
\ee

Gleichung \eqref{eq:normalglbox} hat eine Box, Gleichung \eqref{eq:normalgl} nicht!

Hier ein Verweis auf Bild \vref{fig:bildlabel}.

Hinweis: Hier wird ein Makro zur Einbindung des Bildes benutzt (siehe file 
\texttt{macros.tex}).\index{Bild einbinden!mit Makro}

\figscale{hcos}{Bildunterschrift bla bla bla bla bla bla bla bla bla bla bla bla bla bla bla bla bla bla bla bla bla bla bla bla bla bla bla bla bla bla}{bildlabel}{ht}{.7}
% Hinweis: Skalierung (.7) ist bezogen auf nutzbare Seitenbreite

Und auf die Tabelle \vref{tab:res2wav}. 

\btline{ht}{1.2}
\btab{|c|l|c|c|c|}
 \hline
Nr.& Datei& Dauer in s& \mc{1}{c|}{$G_p/$ dB}& \mc{1}{c|}{$\mbox{seg}G_p/$ dB}\\
 \hline
 \hline
\mc{5}{|l|}{\sf kurze Sinust�ne 14~kHz bis 20~kHz in Halbtonschritten}\\
 \hline
1& \C{38\_2\_24.wav}& 20.06& 29.48& 26.52\\
  \hline
\etab
\et{Liste der wav-Dateien mit Pr�diktionsergebnissen}{res2wav}

\section{Abschnitt zwei}

bla bla bla bla bla bla bla bla bla bla bla bla bla bla bla
bla bla bla bla bla bla bla bla bla bla bla bla bla bla bla
bla bla bla bla bla bla bla bla bla bla bla bla bla bla bla

bla bla bla bla bla bla bla bla bla bla bla bla bla bla bla
bla bla bla bla bla bla bla bla bla bla bla bla bla bla bla
bla bla bla bla bla bla bla bla bla bla bla bla bla bla bla
bla bla bla bla bla bla bla bla bla bla bla bla bla bla bla
bla bla bla bla bla bla bla bla bla bla bla bla bla bla bla
bla bla bla bla bla bla bla bla bla bla bla bla bla bla bla
bla bla bla bla bla bla bla bla bla bla bla bla bla bla bla

bla bla bla bla bla bla bla bla bla bla bla bla bla bla bla
bla bla bla bla bla bla bla bla bla bla bla bla bla bla bla
bla bla bla bla bla bla bla bla bla bla bla bla bla bla bla
bla bla bla bla bla bla bla bla bla bla bla bla bla bla bla
bla bla bla bla bla bla bla bla bla bla bla bla bla bla bla
bla bla bla bla bla bla bla bla bla bla bla bla bla bla bla
bla bla bla bla bla bla bla bla bla bla bla bla bla bla bla

bla bla bla bla bla bla bla bla bla bla bla bla bla bla bla
bla bla bla bla bla bla bla bla bla bla bla bla bla bla bla
bla bla bla bla bla bla bla bla bla bla bla bla bla bla bla
bla bla bla bla bla bla bla bla bla bla bla bla bla bla bla

Listing \ref{lst:bsp} zeigt ein Beispiel f�r das Einbinden einer Datei.
\lstinputlisting[float=htp,
								caption={Beispiel-Listing in eigener Umgebung.}, 
								label=lst:bsp,
								numbers=left, stepnumber=1,
									%firstline=7, lastline=38, 
									stepnumber=1]{sourcecodes/fftdb.m}

Variable \lstinline!e_buff! enth�lt... 

Funktion \lstinline!function fftdb! macht dies und das und jenes


\section{Subfigures}

Beispiel f�r Subfigures

\begin{figure}[ht]
	\centering
	\subfigure[Dummy 1]{
		\includegraphics[width=0.35\columnwidth]{figures/dummy}
		\label{fig:dummy1}
	}
	\qquad
	\subfigure[Dummy 2]{
		\includegraphics[width=0.35\columnwidth]{figures/dummy}
		\label{fig:dummy2}
	}
	\caption{Zwei Dummies in einer Figure}
	\label{fig:dummy12}
\end{figure}

Bild \ref{fig:dummy1} und \ref{fig:dummy1} sind in \ref{fig:dummy12} zu sehen.

\cleardoublepage
\chapter{Zweites Kapitel}
bla bla bla

bla bla bla

\section{Abschnitt zwei-eins}

bla bla bla bla bla bla bla bla bla bla bla bla bla bla bla
bla bla bla bla bla bla bla bla bla bla bla bla bla bla bla
bla bla bla bla bla bla bla bla bla bla bla bla bla bla bla

bla bla bla bla bla bla bla bla bla bla bla bla bla bla bla
bla bla bla bla bla bla bla bla bla bla bla bla bla bla bla
bla bla bla bla bla bla bla bla bla bla bla bla bla bla bla
bla bla bla bla bla bla bla bla bla bla bla bla bla bla bla
bla bla bla bla bla bla bla bla bla bla bla bla bla bla bla
bla bla bla bla bla bla bla bla bla bla bla bla bla bla bla
bla bla bla bla bla bla bla bla bla bla bla bla bla bla bla

bla bla bla bla bla bla bla bla bla bla bla bla bla bla bla
bla bla bla bla bla bla bla bla bla bla bla bla bla bla bla
bla bla bla bla bla bla bla bla bla bla bla bla bla bla bla
bla bla bla bla bla bla bla bla bla bla bla bla bla bla bla
bla bla bla bla bla bla bla bla bla bla bla bla bla bla bla
bla bla bla bla bla bla bla bla bla bla bla bla bla bla bla
bla bla bla bla bla bla bla bla bla bla bla bla bla bla bla

bla bla bla bla bla bla bla bla bla bla bla bla bla bla bla
bla bla bla bla bla bla bla bla bla bla bla bla bla bla bla
bla bla bla bla bla bla bla bla bla bla bla bla bla bla bla
bla bla bla bla bla bla bla bla bla bla bla bla bla bla bla

\section{Abschnitt zwei-zwei}

bla bla bla bla bla bla bla bla bla bla bla bla bla bla bla
bla bla bla bla bla bla bla bla bla bla bla bla bla bla bla
bla bla bla bla bla bla bla bla bla bla bla bla bla bla bla

bla bla bla bla bla bla bla bla bla bla bla bla bla bla bla
bla bla bla bla bla bla bla bla bla bla bla bla bla bla bla
bla bla bla bla bla bla bla bla bla bla bla bla bla bla bla
bla bla bla bla bla bla bla bla bla bla bla bla bla bla bla
bla bla bla bla bla bla bla bla bla bla bla bla bla bla bla
bla bla bla bla bla bla bla bla bla bla bla bla bla bla bla
bla bla bla bla bla bla bla bla bla bla bla bla bla bla bla

bla bla bla bla bla bla bla bla bla bla bla bla bla bla bla
bla bla bla bla bla bla bla bla bla bla bla bla bla bla bla
bla bla bla bla bla bla bla bla bla bla bla bla bla bla bla
bla bla bla bla bla bla bla bla bla bla bla bla bla bla bla
bla bla bla bla bla bla bla bla bla bla bla bla bla bla bla
bla bla bla bla bla bla bla bla bla bla bla bla bla bla bla
bla bla bla bla bla bla bla bla bla bla bla bla bla bla bla

bla bla bla bla bla bla bla bla bla bla bla bla bla bla bla
bla bla bla bla bla bla bla bla bla bla bla bla bla bla bla
bla bla bla bla bla bla bla bla bla bla bla bla bla bla bla
bla bla bla bla bla bla bla bla bla bla bla bla bla bla bla

\section{Abschnitt zwei-drei}

bla bla bla bla bla bla bla bla bla bla bla bla bla bla bla
bla bla bla bla bla bla bla bla bla bla bla bla bla bla bla
bla bla bla bla bla bla bla bla bla bla bla bla bla bla bla

bla bla bla bla bla bla bla bla bla bla bla bla bla bla bla
bla bla bla bla bla bla bla bla bla bla bla bla bla bla bla
bla bla bla bla bla bla bla bla bla bla bla bla bla bla bla
bla bla bla bla bla bla bla bla bla bla bla bla bla bla bla
bla bla bla bla bla bla bla bla bla bla bla bla bla bla bla
bla bla bla bla bla bla bla bla bla bla bla bla bla bla bla
bla bla bla bla bla bla bla bla bla bla bla bla bla bla bla

bla bla bla bla bla bla bla bla bla bla bla bla bla bla bla
bla bla bla bla bla bla bla bla bla bla bla bla bla bla bla
bla bla bla bla bla bla bla bla bla bla bla bla bla bla bla
bla bla bla bla bla bla bla bla bla bla bla bla bla bla bla
bla bla bla bla bla bla bla bla bla bla bla bla bla bla bla
bla bla bla bla bla bla bla bla bla bla bla bla bla bla bla
bla bla bla bla bla bla bla bla bla bla bla bla bla bla bla

bla bla bla bla bla bla bla bla bla bla bla bla bla bla bla
bla bla bla bla bla bla bla bla bla bla bla bla bla bla bla
bla bla bla bla bla bla bla bla bla bla bla bla bla bla bla
bla bla bla bla bla bla bla bla bla bla bla bla bla bla bla

\section{Abschnitt zwei-vier}

bla bla bla bla bla bla bla bla bla bla bla bla bla bla bla
bla bla bla bla bla bla bla bla bla bla bla bla bla bla bla
bla bla bla bla bla bla bla bla bla bla bla bla bla bla bla

bla bla bla bla bla bla bla bla bla bla bla bla bla bla bla
bla bla bla bla bla bla bla bla bla bla bla bla bla bla bla
bla bla bla bla bla bla bla bla bla bla bla bla bla bla bla
bla bla bla bla bla bla bla bla bla bla bla bla bla bla bla
bla bla bla bla bla bla bla bla bla bla bla bla bla bla bla
bla bla bla bla bla bla bla bla bla bla bla bla bla bla bla
bla bla bla bla bla bla bla bla bla bla bla bla bla bla bla

bla bla bla bla bla bla bla bla bla bla bla bla bla bla bla
bla bla bla bla bla bla bla bla bla bla bla bla bla bla bla
bla bla bla bla bla bla bla bla bla bla bla bla bla bla bla
bla bla bla bla bla bla bla bla bla bla bla bla bla bla bla
bla bla bla bla bla bla bla bla bla bla bla bla bla bla bla
bla bla bla bla bla bla bla bla bla bla bla bla bla bla bla
bla bla bla bla bla bla bla bla bla bla bla bla bla bla bla

bla bla bla bla bla bla bla bla bla bla bla bla bla bla bla
bla bla bla bla bla bla bla bla bla bla bla bla bla bla bla
bla bla bla bla bla bla bla bla bla bla bla bla bla bla bla
bla bla bla bla bla bla bla bla bla bla bla bla bla bla bla



\cleardoublepage
\chapter{Drittes Kapitel}


Dies ist ein Text, der nur ein Beispiel darstellt. Dies ist ein Text, der nur ein 
Beispiel darstellt. Dies ist ein Text, der nur ein Beispiel darstellt.

Ein Zitat: \cite{fli:mult}.

Hier ein Verweis auf einen anderen Abschnitt: \ref{sec:abschnittlabel}.

Hier ein Verweis auf Abbildung \ref{fig:bildlabel_1}.

\index{Bild einbinden!ohne Makro}

\begin{figure}[ht]
	\centerline{\includegraphics[scale=0.5]{figures/dummy}}
	\caption{Eine Dummy Abbildung. Eine lange Bildunterschrift. Eine lange Bildunterschrift. Eine lange Bildunterschrift. Eine lange Bildunterschrift. Eine lange Bildunterschrift.}  
	\label{fig:bildlabel_1}
\end{figure}

\begin{figure}[ht]
	\centerline{\includegraphics[width=1.0\columnwidth]{figures/dummy}}
	\caption{Noch eine Dummy Abbildung.}    
	\label{fig:bildlabel_2}
\end{figure}

Dies ist ein Text, der nur ein Beispiel darstellt. Dies ist ein Text, der nur ein 
Beispiel darstellt. Dies ist ein Text, der nur ein Beispiel darstellt.
Dies ist ein Text, der nur ein Beispiel darstellt. Dies ist ein Text, der nur ein 
Beispiel darstellt. Dies ist ein Text, der nur ein Beispiel darstellt.


%\clearpage
%\pagebreak

\section{Abschnitt eins}

\label{sec:abschnittlabel}

Dies ist ein Text, der nur ein Beispiel darstellt. Dies ist ein Text, der nur ein 
Beispiel darstellt. Dies ist ein Text, der nur ein Beispiel darstellt.
Dies ist ein Text, der nur ein Beispiel darstellt. Dies ist ein Text, der nur ein 
Beispiel darstellt. Dies ist ein Text, der nur ein Beispiel darstellt.

\subsection{Unterabschnitt eins-zwei}

Dies ist ein Text, der nur ein Beispiel darstellt. Dies ist ein Text, der nur ein 
Beispiel darstellt. Dies ist ein Text, der nur ein Beispiel darstellt.
Dies ist ein Text, der nur ein Beispiel darstellt. Dies ist ein Text, der nur ein 
Beispiel darstellt. Dies ist ein Text, der nur ein Beispiel darstellt.


\section{Impulsformer - Wurzel-Cosinus-Rolloff Filter}

F�r die Nachrichten�bertragung werden Sende- und Empfangsfilter ben�tigt, die der 
Matched-Filter Bedingung gen�gen. Die Faltung des Sende- mit dem Empfangsfilter sollte 
ferner eine Cosinus-Rolloff-Charakteristik aufweisen. F�r die Sende- und Empfangsfilter 
werden daher Filter mit Wurzel-Cosinus-Rolloff-Charakteristik eingesetzt.

Die zeitkontinuierliche Version eines Cosinus-Rolloff Filters hat den folgenden 
Frequenzgang

\begin{equation}
H_{CR0}(j\omega) = 
\begin{cases}
1 & \mbox{f�r} \; \frac{\mid \omega \mid}{\omega_c} \leq 1-r \\
\frac{1}{2} \biggl(1 + \cos\biggl[ \frac{\pi}{2r}( \frac{\omega}{\omega_c} - 
(1-r))\biggr]  \biggl)
& \mbox{f�r} \; 1-r \leq \frac{\mid \omega \mid}{w_c} \leq 1+r \\
0 & \mbox{f�r} \; \frac{\mid \omega \mid}{\omega_c} \geq 1+r 
\end{cases}
\label{eq:H_CRO}
\end{equation}



\section{pdf features}

Link auf E-Mail:
\href{mailto:dafx2002@unibw-hamburg.de}{dafx2002@unibw-hamburg.de}

Link auf webpage: \\
DAFx main page: \href{http://www.dafx.de}{http://www.dafx.de}

Grafik mit Link auf webpage:
\begin{figure}[ht]
  \centerline{
	\href{http://www.dafx.de}{\includegraphics[width=1.0\columnwidth]{figures/dummy}}
	}
	\caption{Eine Dummy Abbildung mit Link.}    
	\label{fig:bildlabel_3}
\end{figure}



\clearpage

\section{Systemeigenschaften}

Im folgenden sind die verwendeten Parameter mit Zahlenwerten zusammengefasst:
\begin{itemize}
\item Symbol�bertragungsrate $f_T = 1/T = 2 f_N = 24 \; \mbox{kHz}$ \\
\item Bit�bertragungsrate bei verwendeter QPSK-Modulation $f_B = 2 f_T = 48 \; 
\mbox{kBit/s}$ \\
\item Symboldauer $T = \frac{1}{2 f_N} = 41.67 \; \mu$s \\
\item Nyquistfrequenz $f_N = 12 \; \mbox{kHz}$ \\
\item Rolloff-Faktor $r = 0.5$ \\
\item Parameter $M$ gibt den Faktor der �berabtastung an, $M = 2,3,4...$. \\ Bei 
verwendeter vierfacher �berabtastung $M = 4$. \\
\item Bandbreite im �bertragungsband (siehe Abbildung \ref{fig:bildlabel_3}) 
\begin{align}
B & = 2 \; (1+r) f_N \\
& = (1+r) f_T \\
& = 36 \; \mbox{kHz}
\end{align}
\item mindestens notwendige Abtastfrequenz nach Abw�rtsmischung ins Basisband 
\begin{align}
f_A & = 2 \; [2 \; f_N] \geq 2 \; [(1+r) \; f_N] \\
& = \frac{2}{T} = 2 f_T \\
& = 48 \; \mbox{kHz}
\end{align}
\item Abtastfrequenz bei verwendeter vierfacher �berabtastung nach Abw�rtsmischung ins 
Basisband 
\begin{align}
f_A & = 4 \; [2 \;f_N] \geq 4 \; [(1+r)\; f_N] \\
& = \frac{4}{T} = 4 f_T \\
& = 96 \; \mbox{kHz}
\end{align}
\item Tr�gerfrequenz $f_0 = 800 \; \mbox{MHz} \; / \; 2 \; \mbox{GHz}$ \\
\end{itemize}



%############# appendix #################
\cleardoublepage
\renewcommand{\chaptername}{\appendixname} %if name appears in header

\begin{appendix}

\chapter{Anhang 1}
bla bla bla

bla bla bla



\cleardoublepage
\chapter{Anhang 2}
bla bla bla

bla bla bla



\end{appendix}

\cleardoublepage
\phantomsection
\addcontentsline{toc}{chapter}{\protect\numberline{Abk�rzungsverzeichnis}}
\chapter*{Abk\"urzungsverzeichnis}
\label{sec:abkuerzungen}
\markboth{\MakeUppercase{Abk\"urzungsverzeichnis}}{\MakeUppercase{Abk\"urzungsverzeichnis}} 
\begin{longtable}[t]{ll}
\arrayrulecolor{hsugrau}
\textbf{\textcolor{hsurot}{A}} &\\ \hline 
ANT & Allgemeine Nachrichtentechnik \\

 
  &\\ \textbf{\textcolor{hsurot}{H}} &\\  \hline 
HSU-HH & Helmut-Schmidt-Universit\"at/ Universit\"at der Bundeswehr Hamburg\\

\end{longtable}

%Softwareverzeichnis
\cleardoublepage
\phantomsection
\addcontentsline{toc}{chapter}{\protect\numberline{Softwareverzeichnis}}
\chapter*{Softwareverzeichnis}
\markboth{\MakeUppercase{Softwareverzeichnis}}{\MakeUppercase{Softwareverzeichnis}} 
\label{Softwareverzeichnis}
\begin{table}[htb]
	\centering
		\begin{tabular}{|l|l||l|l|}
		\hline
			\bfseries{Name} & \bfseries{Version} & \bfseries{URL} & \bfseries{Bemerkung}\\
		\hline
		\hline
			Java Development Kit & 1.5 U6 & http://java.sun.com& Java SDK\\
		\hline
			JDBC & --- & s.o. & Datenbank API\\
		\hline
			Java3D & --- & s.o. & 3D-API \\
		\hline
			Java Advanced Imaging & & s.o. & Bildbearbeitungs API auf \\& & &Grundlage von Systembibliotheken\\
		\hline
			JAMA & --- & s.o.  & Matrixpaket\\
		\hline
			MySQL & 4.1& www.mysql.de & Datenbank\\
		\hline
			MySQL & 5& s.o.& Datenbank\\
		\hline
			MySQL JDBC Treiber & --- & s.o. & Datenbanktreiber \\
		\hline
		\end{tabular}
	\caption{Verwendete Software.}
	\label{tab:VerwendeteSoftware}
\end{table}

blablabla

\cleardoublepage
%########### Bibliography ##############
\clearpage % to get correct page no. in TOC
\phantomsection
\addcontentsline{toc}{chapter}{\protect\numberline{\bibname}}

\renewcommand{\chaptername}{} %if name appears in header

\renewcommand{\baselinestretch}{1}
\normalsize

% if enabling this command all references of the bib file are listed
% otherwise only the referenced books are listed in Bibliography section
%\nocite{*} 

%\bibliographystyle{unsrt}
%\bibliographystyle{alpha}
%\bibliographystyle{ALPHADIN}

\bibliographystyle{newalpha} %neuer Stil mit makebst erzeugt
% newalpha.dbj wurde durch "latex makebst" erzeugt
% und enth�lt die Formatierungen (file kann per Hand editiert werden)
% mit "latex newalpha.dbj" wird das .bst file erzeugt

	
% In literatur.bib sind die eigentlichen Literaturangaben enthalten
\bibliography{literatur} % requires file literatur.bib


\ifmakeindex
\cleardoublepage
\phantomsection
\addcontentsline{toc}{chapter}{\protect\numberline{Index}}
% Stichwortverzeichnis endgueltig anzeigen
\printindex
\fi

\end{document}
